\documentclass[11pt,a4paper]{article}

\usepackage[T1]{fontenc}
\usepackage[utf8]{inputenc}
\usepackage{lmodern}
\usepackage{microtype}
\usepackage[margin=25mm]{geometry}
\usepackage{amsmath,amssymb,amsthm,mathtools}
\usepackage{booktabs,longtable,array}
\usepackage{enumitem}
\usepackage{pdflscape}
\usepackage{xcolor}
\usepackage{url}
\usepackage[hidelinks]{hyperref}
\usepackage[nameinlink,noabbrev]{cleveref}

\hypersetup{
  pdftitle={Supplementary proofs and verification dossier for The exact saturated 6- and 7-Sperner numbers},
  pdfauthor={Jiazhi Chen},
  pdfsubject={Arbitrary-finite blocker reductions, finite classifications, and Lean 4 verification}
}

\newtheorem{theorem}{Theorem}[section]
\newtheorem{lemma}[theorem]{Lemma}
\newtheorem{proposition}[theorem]{Proposition}
\newtheorem{corollary}[theorem]{Corollary}
\theoremstyle{definition}
\newtheorem{definition}[theorem]{Definition}
\newtheorem{remark}[theorem]{Remark}

\newcommand{\powerset}{2^{}}
\newcommand{\blocker}{B}
\newcommand{\sat}{\operatorname{sat}}
\newcommand{\family}{\mathcal F}
\newcommand{\minrows}{\operatorname{min}}
\newcommand{\code}[1]{\texttt{#1}}
\newcommand{\proved}{\textsc{Proved}}
\newcommand{\computed}{\textsc{Computed}}
\newcommand{\formalized}{\textsc{Formalized}}
\renewcommand{\thesection}{S\arabic{section}}
\renewcommand{\thesubsection}{\thesection.\arabic{subsection}}
\makeatletter
\renewcommand*\l@section{\@dottedtocline{1}{1.5em}{4em}}
\renewcommand*\l@subsection{\@dottedtocline{2}{2.5em}{4.5em}}
\makeatother

\title{Supplementary proofs and verification dossier for\\
\emph{The exact saturated 6- and 7-Sperner numbers}}
\author{Jiazhi Chen\\
\small Hangzhou No.2 High School of Zhejiang Province\\
\small \href{mailto:lailai0x394@gmail.com}{lailai0x394@gmail.com}\quad
\href{https://orcid.org/0009-0008-9790-9881}{ORCID 0009-0008-9790-9881}}
\date{July 30, 2026}

\begin{document}
\maketitle

\begin{abstract}
This supplement records the arbitrary-finite reductions and the exact finite
classification interfaces used in the proof that
$\sat(6)=30$ and $\sat(7)=55$.  Its purpose is auditability.  In particular,
it explains why the finite searches do not impose a bound on the original
ground set, states every residual search domain, records independent replay
counts, and maps the natural-language statements to their Lean 4
declarations.  Inactive points, parallel actual points, repeated incidence
patterns, nonuniform rows, and empty optional incidence classes are retained
throughout.  For the larger classifications, the named Lean declarations are
the primary proof objects; the prose below records their mathematical
interfaces and an audit path rather than reproducing every kernel reduction
line by line.
\end{abstract}

\tableofcontents

\section{Evidence labels and common blocker lemmas}
\label{sec:evidence}

All hypergraphs below are finite clutters on an arbitrary finite set $U$.
For a clutter $\mathcal H$, $\blocker(\mathcal H)$ is the clutter of its
inclusion-minimal transversals.  Three evidence labels are used strictly:

\begin{description}[style=nextline,leftmargin=29mm,font=\normalfont]
\item[\proved]
an argument written out from the displayed definitions and lemmas, or an
independently checked proof certificate for its stated finite encoding;
\item[\computed]
an output of a finite executable enumeration, independently replayed but not
by itself a theorem about arbitrary finite clutters;
\item[\formalized]
a proposition accepted by the Lean kernel from the checked sources.
\end{description}

The C++ and Python programs are discovery and replay tools.  The final exact
claims do not trust their output directly: the arbitrary-finite reductions
are formalized, and the residual finite classifications are represented by
kernel-checked Lean proofs.  The executable replays remain useful because
they provide a second implementation and make the classification domains
easy to inspect.  A compressed numerical summary is not assigned the
\proved\ label merely because a replay accepts it; its arbitrary-finite
conclusion receives the \formalized\ label only through the named Lean
soundness and realization declarations.

\begin{lemma}[Private rows]
\label{lem:private}
If $T\in\blocker(\mathcal H)$ and $x\in T$, then some
$H_x\in\mathcal H$ satisfies $T\cap H_x=\{x\}$.  Distinct points of $T$
have distinct private rows, and hence $|T|\leq|\mathcal H|$.
\end{lemma}

\begin{proof}
The set $T\setminus\{x\}$ is not a transversal, so it misses a row $H_x$.
Since $T$ is a transversal, $T\cap H_x=\{x\}$.  One row cannot be private
for two different points.
\end{proof}

\begin{lemma}[Residual blockers]
\label{lem:residual}
Fix $x\in U$.  Delete $x$ from every row of $\mathcal H$ and retain only
the inclusion-minimal residual rows, obtaining $\mathcal R_x$.  Then
\[
  \blocker(\mathcal R_x)
  =\{T\in\blocker(\mathcal H):x\notin T\}.
\]
If every residual row has size at least two, the right-hand side has at least
two members.
\end{lemma}

\begin{proof}
A set avoiding $x$ meets every row of $\mathcal H$ precisely when it meets
every residual row.  Minimality is unchanged after redundant residual rows
are deleted.  If $\blocker(\mathcal R_x)=\{Q\}$, blocker involution makes
$\mathcal R_x$ the family of singleton rows $\{\{q\}:q\in Q\}$, contrary
to the assumed row-size bound.
\end{proof}

\begin{lemma}[Injection through a point]
\label{lem:injection}
Let $\mathcal H_0$ be the rows of $\mathcal H$ avoiding $x$.  For each
$T\in\blocker(\mathcal H_0)$, use $T$ if it already meets every row of
$\mathcal H$ and otherwise minimize $T\cup\{x\}$ while retaining $T$.
This gives an injection
\[
  \blocker(\mathcal H_0)\hookrightarrow\blocker(\mathcal H).
\]
Every pair blocker of $\mathcal H_0$ must acquire $x$ if all blockers of
$\mathcal H$ have size at least three.
\end{lemma}

\begin{proof}
Rows in $\mathcal H_0$ supply private witnesses for the points of $T$.  In
the second case, a row missed by $T$ supplies a private witness for $x$.
The intersection with $U\setminus\{x\}$ recovers $T$, so the map is
injective.
\end{proof}

\begin{lemma}[Three-row formula]
\label{lem:three-row}
Let $\mathcal K=\{K_1,K_2,K_3\}$ have row sizes at least three and
$\tau(\mathcal K)\geq2$.  Write $p,q,r$ for the multiplicities of the
three private incidence classes and $x,y,z$ for the three pairwise classes.
Then
\[
 |\blocker(\mathcal K)|
 =xy+xz+yz+xr+yq+zp+pqr.                 \tag{S1.1}
\]
The first six terms count pair blockers and $pqr$ counts triple blockers.
In particular $|\blocker(\mathcal K)|\geq7$, and equality permits at most
one triple blocker.
\end{lemma}

\begin{proof}
The triple intersection is empty.  A minimal transversal either selects two
compatible pairwise/private incidence classes or one point from every
private class; private rows exclude every other type.  This gives (S1.1).
The six row-size inequalities then give the lower bound.  If $pqr\geq2$,
classification by the number of nonzero variables among $x,y,z$ gives at
least nine terms, so equality seven has $pqr\leq1$.
\end{proof}

\section{The arbitrary-finite lower bound \texorpdfstring{$m(3,3)\geq14$}{m(3,3)>=14}}
\label{sec:m33-lower}

Let $\mathcal S=\blocker(\mathcal C)$ and
$\mathcal C=\blocker(\mathcal S)$, and suppose every row on both sides has
size at least three.  The preceding private-row and residual arguments first
give $|\mathcal S|+|\mathcal C|\geq13$.  If equality held, blocker duality
would allow us to orient the pair so that
\[
 (|\mathcal S|,|\mathcal C|)\in\{(4,9),(5,8),(6,7)\}.       \tag{S2.1}
\]
We record the sharpened bounds that exclude these three cases.

The paragraphs below give the mathematical reduction used to organize these
branches.  They are not intended as a line-by-line replacement for the larger
case analyses in the formal development.  The arbitrary-finite proof object
for the assembled conclusion is
\path{MiddleLower.mutual_blocker_total_card_ge_fourteen}; its hypotheses retain
the original finite type, both blocker equalities, clutter minimality, and the
two row-size bounds.

\begin{lemma}[Four rows]
\label{lem:four-nineteen}
If a four-row clutter $\mathcal H$ has row sizes and transversal number at
least three, then $|\blocker(\mathcal H)|\geq19$.
\end{lemma}

\begin{proof}
The row-intersection graph has matching number at most one.  With two
isolated vertices, blocker factorization gives at least $3^3=27$ blockers.
With one isolated row, \cref{lem:three-row} gives at least $3\cdot7=21$.
With no isolated row the graph is a three-leaf star.  If $C$ is its center,
$L_i$ its leaves, $X_i=C\cap L_i$, $Y_i=L_i\setminus C$, and
$Z=C\setminus\bigcup_iL_i$, put $x_i=|X_i|$, $y_i=|Y_i|$,
$z=|Z|$, and $n_i=x_i+y_i\geq3$.  Directly listing the minimal
transversals gives
\[
 |\blocker(\mathcal H)|
 =n_1n_2n_3+(z-1)y_1y_2y_3.
\]
For $z\geq1$ this is at least 27.  For $z=0$ it is at least
$n_1n_2n_3-(n_1-1)(n_2-1)(n_3-1)\geq19$.
\end{proof}

\begin{lemma}[Five rows]
\label{lem:five-nine}
If a five-row clutter $\mathcal H$ has row sizes and transversal number at
least three, then $|\blocker(\mathcal H)|\geq9$.
\end{lemma}

\begin{proof}
For each point $e$, \cref{lem:residual} and blocker involution show that at
least four blockers avoid $e$: one or two residual blockers force singleton
or too many residual rows, while three invoke \cref{lem:three-row} and force
at least seven residual rows.  In the only tight row-intersection branch,
two actual matching edges share $e$.  The standard private-row construction
gives at least five blockers containing $e$ (six when the two remaining
endpoints are parallel actual points).  The four avoiding blockers are
distinct from them.  All other intersection-graph branches already give
at least nine by factorization or independent choices.
\end{proof}

\begin{lemma}[Six rows]
\label{lem:six-eight}
If a six-row clutter $\mathcal H$ has row sizes and transversal number at
least three, then $|\blocker(\mathcal H)|\geq8$.
\end{lemma}

\begin{proof}
A point in at least four rows gives at least nine blockers by private choices
from the avoiding rows and \cref{lem:residual}.  If a point $x$ has degree
three, \cref{lem:injection} applied to the three avoiding rows gives seven
blockers.  When equality holds, \cref{lem:three-row} says at most one of
these can avoid $x$, while \cref{lem:residual} gives two avoiding blockers;
one is new, giving eight.

It remains that all degrees are at most two.  Encode every actual point as a
loop or nonloop edge on the six row labels; parallel actual points are not
identified.  The matching-number zero, one, and two cases give at least nine
minimal edge covers.  For a perfect matching $M$, delete its three actual
points and minimize the residual rows.  There are at least four blockers
avoiding $M$.  If the original blocker had at most seven rows, at most two
additional rows $N_j$ would meet $M$.  Write
$I_j=N_j\cap M$ and $P_j=N_j\setminus M$.  For $e\in M$, the blocker of
\[
 \minrows\bigl(\mathcal D\cup\{P_j:e\notin I_j\}\bigr)
\]
injects into rows whose intersection with $M$ is $\{e\}$.  With at most one
residual row it has at least three members; with two residual rows of size at
least two it has at least two.  Summing over the three points of $M$ gives at
least seven rows in the double blocker, contradicting that it is the
six-row clutter $\mathcal H$.  Thus the original blocker has at least eight
rows.
\end{proof}

\begin{theorem}
\label{thm:m33-lower}
Every mutual-blocker pair with minimum row size three on both sides has total
cardinality at least fourteen.
\end{theorem}

\begin{proof}
The three cases in (S2.1) contradict, respectively,
\cref{lem:four-nineteen,lem:five-nine,lem:six-eight}.
\end{proof}

The argument is arbitrary-finite.  Inactive points never enter a minimal
transversal; the loop-multigraph representation retains every parallel
actual point; and minimizing residual rows permits repeated supports and
nonuniform rows.

\section{Equality fourteen, the Fano plane, and exclusion of total fifteen}
\label{sec:m33-classification}

The seven lines of the Fano plane form a self-blocking clutter, so
\cref{thm:m33-lower} is sharp and $m(3,3)=14$.  What is needed globally is
the stronger equality classification and the exclusion of one unit of
slack.

For total fourteen, the asymmetric ordered splits are
$(5,9)$ and $(6,8)$ and their duals.  The five-row case has no point of degree
three or more: such a point produces at least nine blockers containing it
and another avoiding it.  At maximum degree two, choose an inclusion-minimal
actual loop-multigraph core in which every row degree is at least three.
Every support touches a degree-three vertex.  If $l$ is the number of such
vertices, $4(5-l)\leq3l$, so $l\in\{3,4,5\}$.  Exact enumeration of all
multiplicities on the 15 loop/edge supports gives minimum numbers of actual
minimal edge covers $19,16,15$, respectively.  Thus $(5,9)$ is impossible.

For $(6,8)$, degree at least four gives nine blockers.  A degree-three point
leaves a three-row kernel.  Formula (S1.1), the residual blocker lemma, and
the required blocker sizes exclude all low-count kernels; the unique
borderline vector is
\[
 (p,q,r;x,y,z)=(1,1,1;1,1,1),
\]
with six pair blockers and one triple blocker, and its two-row residual would
have either a singleton or at least nine cross-pairs.  At maximum degree two,
the exact six-row loop-multigraph classification has at least nine minimal
edge covers.  Hence $(6,8)$ is also impossible.

In the remaining $(7,7)$ case, a three-point row is forced by the Bernoulli
cover inequality.  Degree and residual reductions leave 41 four-row support
kernels.  Only the six-edge incidence clutter of $K_4$ has at most three
pair blockers; injecting those three pairs through the removed point turns
its four stars and three matchings into the seven Fano lines.  Blocker
involution fixes the opposite side.  Consequently equality fourteen is,
after deleting inactive points and relabelling active points, precisely the
Fano line clutter paired with itself.

For total fifteen, blocker duality leaves splits $(5,10)$, $(6,9)$, and
$(7,8)$.  The five-row loop-multigraph bound above is at least fifteen, so it
excludes $(5,10)$.  For $(6,9)$, high degrees are excluded by injections and
residual blockers; the only degree-three three-row kernel satisfying
$b_2\leq5$ and $b_2+b_3\leq9$ is
\[
 (p,q,r;x,y,z)=(1,2,2;1,1,0),\qquad(b_2,b_3)=(5,4).
\]
The five pair images and four triple images exhaust the nine blockers and
force a two-point row in the double blocker, a contradiction.  At maximum
degree two, actual minimal edge covers already give at least ten blockers.

Finally consider $(7,8)$.  Degree bounds give maximum degree three on both
sides.  The Bernoulli cover inequality forces a three-point row.  If it lies
on the eight-row side, deleting a degree-three point leaves the unique
$K_4$ kernel; its seven images are Fano, and the eighth row is disjoint from
the Fano support, giving at least $7\cdot3$ double-blocker rows.  If it lies
on the seven-row side, the remaining five-row kernel uses the 25 nonempty
supports of rank at most three.  The complete domain contains 490
multiplicity vectors in eight row-permutation classes, all with exactly
three pair blockers.  The 55 possible one-row extensions all fail direct
blocker reconstruction.  Thus total fifteen is impossible.

\begin{theorem}[Middle classification below sixteen]
\label{thm:m33-classification}
One has $m(3,3)=14$; equality is uniquely Fano up to inactive points and
relabeling; and no admissible pair has total cardinality fifteen.
\end{theorem}

The numerical classifications in this section are \computed\ when read as
the C++/Python replay.  The arbitrary-finite conclusions, including the
support and realization interfaces, are also \formalized; their exact Lean
declarations are listed in \cref{sec:crosswalk}.

\section{Canonical slack profiles 52, 53, and 54}
\label{sec:profiles}

The canonical reduction in the main article gives the layerwise baseline
\[
 (|\mathcal A_0|,\ldots,|\mathcal A_6|)
 \geq(1,6,12,14,12,6,1).                \tag{S4.1}
\]
Equality in the middle position is Fano by
\cref{thm:m33-classification}.  The adjacent equality interface says that
each Fano line strictly contains a distinct two-point row in the neighboring
small trace clutter.  The possible adjacent splits are $(6,6)$, $(7,5)$,
and $(8,4)$.  The first lacks the seven required distinct pairs.  For
$(7,5)$ the selected pairs form a seven-edge graph on the Fano points; every
minimal vertex cover has size at least four, and the exact graph argument
gives at least six covers.  For $(8,4)$, the four opposite rows are
pairwise intersecting, have size at least four, and have no singleton
transversal; the four-row trace argument gives at least nine blockers.  All
three splits are impossible.  This excludes total 52.

At total 53, if the middle total remains fourteen, the unique slack lies in
one adjacent layer and the equality layer on the other side is excluded by
the preceding Fano interface.  The only other profile would have middle
total fifteen, excluded by \cref{thm:m33-classification}.  Hence total 53 is
impossible.

At total 54, distribute two units of slack over the five internal layers.
A fourteen-member middle pair is Fano and forces both adjacent layers above
equality.  A fifteen-member middle pair is impossible.  A larger middle
pair must therefore have total sixteen and leaves both adjacent layers at
equality.  The only profiles are
\[
 P_{\rm F}=(1,6,13,14,13,6,1),\qquad
 P_{16}=(1,6,12,16,12,6,1).             \tag{S4.2}
\]

\section{The Fano profile}
\label{sec:fano-profile}

Let $X$ be the seven active Fano points and $\mathcal L$ its line clutter.
For the lower adjacent pair write $\mathcal F=\mathcal S_2$ and
$\mathcal G=\mathcal C_2$.  Each Fano line strictly contains a distinct pair
from $\mathcal F$.  Up to blocker duality, the possible splits are
$(7,6)$, $(8,5)$, and $(9,4)$.

\subsection{The split \texorpdfstring{$(7,6)$}{(7,6)}}

The seven rows of $\mathcal F$ select one of the three pairs on each Fano
line, so they form a simple graph $Q$ on $X$.  Points outside $X$ occur in no
row and cannot enter a minimal transversal.  This is an exact reduction to
$3^7=2187$ graphs.  The independently replayed blocker histogram among the
1011 graphs whose minimal vertex covers all have size at least four is
\[
 |\blocker(E(Q))|=7:24,\qquad8:903,\qquad9:84.             \tag{S5.1}
\]
No graph has six such covers, excluding $(7,6)$.

\subsection{The split \texorpdfstring{$(8,5)$}{(8,5)}}

Write $\mathcal F=E(Q)\cup\{R\}$,
$A=R\cap X$, $Z=R\setminus X$, and $t=|Z|$.  The clutter condition makes
$A$ independent in $Q$, and \cref{lem:private} gives
\[
 2\leq|A|+t=|R|\leq5.                   \tag{S5.2}
\]
Thus the arbitrary outside set is retained exactly by its multiplicity $t$,
which is proved bounded; no ground-set bound is assumed.  If
$\mathcal V=\blocker(E(Q))$, the blockers of $\mathcal F$ are the minimal
members of
\[
\begin{split}
 \mathcal W(Q,A,Z)={}&\{C\in\mathcal V:C\cap A\ne\varnothing\}\\
 &\cup\{C\cup\{a\}:C\in\mathcal V, C\cap A=\varnothing, a\in A\}\\
 &\cup\{C\cup\{z\}:C\in\mathcal V, C\cap A=\varnothing, z\in Z\}.
\end{split}                                                  \tag{S5.3}
\]
Indeed, minimize a transversal first against $E(Q)$; if its minimal cover
misses $A$, exactly one point of $A\cup Z$ remains.

Equations (S5.2)--(S5.3) give exactly 249576 triples $(Q,A,t)$.  Of these,
8400 have five blockers, but every such case has a three-point blocker.
After imposing minimum blocker size four, 173304 triples remain and the
minimum blocker count is six.  Hence $(8,5)$ is impossible.

\subsection{The split \texorpdfstring{$(9,4)$}{(9,4)}}

We use the following purely combinatorial bound.

\begin{lemma}
\label{lem:four-ten}
If $\mathcal D=\{D_1,D_2,D_3,D_4\}$ is a clutter whose rows have size at
least four, every two rows intersect, and $\tau(\mathcal D)\geq2$, then
$|\blocker(\mathcal D)|\geq10$.
\end{lemma}

\begin{proof}
Choose a smallest row $E=D_1$, $n=|E|$, put
$O_j=E\setminus D_j$, and for $x\in E$ set
$w_x=\min_{j:x\in O_j}|O_j|$ and $W=\sum_xw_x$.  Private rows give $W$
blockers meeting $E$ exactly in one point, while an assignment of every
$x$ to a minimizing omission gives
\[
 W\geq\left\lceil\frac{n^2}{3}\right\rceil.               \tag{S5.4}
\]
The blocker of the trace clutter
$\{E,E\cap D_2,E\cap D_3,E\cap D_4\}$ supplies further blockers lying
inside $E$, disjoint from the $W$ classes.  For $n\geq5$, (S5.4) gives nine
and one trace blocker suffices.  For $n=4$, the possible values
$W=6,7,8$ force respectively at least five, three, or two trace blockers;
these follow by blocker involution on one or two trace rows and the omission
partition $2+1+1$.  Thus the total is at least ten.
\end{proof}

Every row of $\mathcal G$ strictly contains a Fano line, so the four rows are
pairwise intersecting, have size at least four, and have transversal number
at least two.  \Cref{lem:four-ten} gives
$|\mathcal F|\geq10$, contradicting the split.  Blocker duality excludes the
upper adjacent orientations, so $P_{\rm F}$ is impossible.

\section{A three-point row and the split \texorpdfstring{$(5,11)$}{(5,11)}}
\label{sec:five-eleven}

Let $\mathcal D$ and $\mathcal E=\blocker(\mathcal D)$ be the two sides of
the middle layer, with
\[
 |\mathcal D|+|\mathcal E|=16,\qquad
 \min_{D\in\mathcal D}|D|,
 \min_{E\in\mathcal E}|E|\geq3.          \tag{S6.1}
\]
For a uniformly random $R\subseteq U$, the events
\[
 D\subseteq R\quad(D\in\mathcal D),\qquad
 E\subseteq U\setminus R\quad(E\in\mathcal E)             \tag{S6.2}
\]
form an exhaustive cover.  An event of the first orientation is disjoint
from an event of the second because every $D$ meets every $E$.  Events on
the same side can overlap.  If all sixteen rows had size at least four, the
union bound would give equality in
\[
 1\leq\sum_D2^{-|D|}+\sum_E2^{-|E|}\leq1.                 \tag{S6.3}
\]
Equality would force every event in (S6.2) to be pairwise disjoint.  Each
side has at least three rows by \cref{lem:private}, whereas two same-side
events have positive-probability intersection.  This contradiction proves
that some middle row has size three.

Now suppose $(|\mathcal D|,|\mathcal E|)=(5,11)$.  No point belongs to four
rows of $\mathcal D$, because it could be extended with one point from the
remaining row to a transversal of size at most two.  If all point degrees
were at most two, the inclusion-minimal actual loop-multigraph core used in
\cref{sec:m33-classification} would already have at least fifteen minimal
edge covers.  Hence some $x$ has degree three.

Let $P,Q$ be the two rows avoiding $x$.  They are disjoint.  Each
$\{x,p,q\}$, $p\in P$, $q\in Q$, is a minimal blocker, so at least nine
blockers contain $x$.  The residual-blocker lemma gives at least two blockers
avoiding $x$.  Equality is forced throughout.  If these two residual
blockers $E_0,E_1$ intersected, a common point together with $x$ would be a
two-point transversal of $\mathcal E$.  Thus they are disjoint and both have
size at least three.  Blocker involution gives
\[
 \mathcal R=\blocker(\{E_0,E_1\})
 =\{\{u,v\}:u\in E_0, v\in E_1\},
\]
so $|\mathcal R|\geq9$.  But $\mathcal R$ was obtained by minimizing only
five residual rows.  This contradiction excludes $(5,11)$ and, by blocker
duality, $(11,5)$.

\section{The split \texorpdfstring{$(6,10)$}{(6,10)}}
\label{sec:six-ten}

Assume $|\mathcal D|=6$ and $|\mathcal E|=10$.  Every point has
$\mathcal D$-degree at most three.  Degrees five and six immediately give a
transversal of size at most two.  At degree four, the two avoiding rows are
disjoint and give at least nine blockers containing the point; the residual
blocker lemma gives two more avoiding it.

\subsection{The degree-three branch}

Let $x$ have degree three and let $\mathcal D_0$ be the three rows avoiding
$x$.  Put $L_j=X_j\setminus\{x\}$ for the three rows containing $x$ and
\[
 \mathcal R=\minrows(\mathcal D_0\cup\{L_1,L_2,L_3\}),
 \qquad
 \mathcal T=\blocker(\mathcal R)
 =\{E\in\mathcal E:x\notin E\}.          \tag{S7.1}
\]
The three-row formula and point-degree bounds on the opposite side imply
that the number $b_2$ of pair blockers of $\mathcal D_0$ is at most six and
$b_2+b_3\leq10$.  Up to row permutation, the only incidence vectors are
\[
\begin{array}{c@{\qquad}c@{\qquad}c}
(1,1,1;1,1,1)&b_2=6&b_3=1,\\
(1,2,2;1,1,0)&b_2=5&b_3=4,\\
(1,1,3;2,0,0)&b_2=6&b_3=3.
\end{array}                                                   \tag{S7.2}
\]
This follows by classifying the number of positive pair classes in
(S1.1): zero gives $pqr\geq27$; one forces the last vector; two force unit
pair multiplicities and the middle vector; three force all six
multiplicities to be one.

The last two kernels are excluded algebraically.  For
$(1,1,3;2,0,0)$ the six pair images use all possible occurrences of $x$.
The three triple blockers and a fourth row $Q$ make up $\mathcal T$.
The row $Q$ must avoid the two common private points; each $q\in Q$ then
creates two distinct two-point blockers, and the remaining three-point class
creates a further blocker, forcing $|\mathcal R|>6$.

For $(1,2,2;1,1,0)$, the five pair images leave at most one triple image
containing $x$.  If none does, the four triples and one row $Q$ force five or
six pair blockers of $\mathcal T$.  If one does, the remaining three triples
form a four-vertex path after deleting their common point, and its three
minimal vertex covers again force $|\mathcal R|=6$.  In either case all rows
of $\mathcal R$ are pairs, so none of the three length-three rows of
$\mathcal D_0$ survives in (S7.1).  Only $L_1,L_2,L_3$ remain available,
giving $|\mathcal R|\leq3$, a contradiction.

The first vector in (S7.2) has realization
\[
 \{\{a,u,v\},\{b,u,w\},\{c,v,w\}\}.      \tag{S7.3}
\]
Its six pair images contain $x$ and its unique triple blocker
$\{a,b,c\}$ avoids $x$.  Hence
\[
 \mathcal T=\{\{a,b,c\},Q_1,Q_2,Q_3\}.  \tag{S7.4}
\]
Every row in (S7.4) has size between three and six.  No point has full
four-row support, so there are fourteen nonempty proper supports.  The
finite domain consists of all $2^{14}-1=16383$ active support masks and all
positive multiplicities satisfying the four row-size bounds.  The exact
results are
\[
\begin{array}{c|c|r}
|\mathcal R|&\text{pair rows in }\mathcal R&\text{vectors}\\
\hline
5&5&6\\
6&5&12\\
6&6&77
\end{array}                                                   \tag{S7.5}
\]
Thus there are 95 vectors in nine row-permutation classes.  For each one,
the semantic check reconstructs $\mathcal R=\blocker(\mathcal T)$ and
exhausts every embedding of (S7.3): if $|\mathcal R|=6$, all three rows in
(S7.3) must survive; if it is five, at least two must survive.  There are no
compatible extensions.

\subsection{The maximum-degree-two branch}

Represent each actual point as a loop or edge on the six row labels, keeping
parallel points distinct.  Choose an inclusion-minimal actual core $J$ in
which every row degree is at least three.  Every support touches a
degree-three vertex.  If $l$ is the number of degree-three vertices, then
\[
 4(6-l)\leq3l,
\]
so $l\in\{4,5,6\}$.  All positive supports touch one of these vertices and
their actual multiplicities are bounded by its exact degree three.  This is
a finite domain on the 21 loop/edge supports, with 562 minimal support
covers.  For multiplicity vector $m$, the exact number of actual minimal edge
covers is
\[
 N(m)=\sum_M\prod_{s\in M}m_s.                              \tag{S7.6}
\]
The independently replayed minima, refined by nonloop-support matching
number $\nu$, are
\[
\begin{array}{c|r|rrrr}
l&\text{vectors}&\nu=0&\nu=1&\nu=2&\nu=3\\
\hline
4&5244&--&--&54&24\\
5&57080&--&195&48&21\\
6&44288&729&171&45&15
\end{array}                                                   \tag{S7.7}
\]
Thus every core has at least fifteen blockers, contradicting ten.  Together
with the degree-three branch this excludes $(6,10)$ and $(10,6)$.

\section{The split \texorpdfstring{$(8,8)$}{(8,8)}}
\label{sec:eight-eight}

Let $|\mathcal D|=|\mathcal E|=8$.  Point degrees at least seven give a
transversal of size at most two; degree six leaves two rows whose private
choices give at least nine blockers.  If a point had degree five, the three
avoiding rows would have at most eight blockers and at most five pair
blockers.  Formula (S1.1) has no such integer solution.  Hence degrees are at
most four.

At degree four, the four avoiding rows have sizes three through eight, at
most eight blockers, and at most four pair blockers.  Their points use the
fourteen nonempty proper supports.  Among all 16383 active support masks and
all bounded multiplicities, the exact domain has
\[
\begin{array}{c|c|r}
|\blocker|&\text{pair blockers}&\text{vectors}\\
\hline
7&3&1\\
8&4&48.
\end{array}                                                   \tag{S8.1}
\]
The semantic extension checker permits every incidence of existing points
consistent with the degree bound and every multiplicity of new singleton-
support points.  It checks 1432 candidates and finds none.  Degree four is
therefore impossible on either side.

By \cref{sec:five-eleven}, some row $T$ has size three.  Put
$T\in\mathcal D$.  It covers all eight rows of $\mathcal E$, and at least two
points of $T$ have $\mathcal E$-degree three.  Delete the three rows through
one such point $x$.  The remaining five-row kernel has sizes three through
eight, at most eight blockers, at most three pair blockers, and supports of
rank at most three.  There are 25 such nonempty supports.  The complete
enumeration of $2^{25}-1=33554431$ active masks gives
\[
\begin{array}{lr}
\text{support-feasible masks}&96148\\
\text{multiplicity vectors}&1890\\
\text{row-permutation classes}&24\\
\text{vectors with 6/7/8 blockers}&20/470/1400.
\end{array}                                                   \tag{S8.2}
\]
Every vector has exactly three pair blockers; all three acquire $x$.  If the
kernel has six, seven, or eight blockers, two, one, or no additional
$x$-free rows remain.  New actual points may have either singleton support
or the shared support of two new rows; all allowed multiplicities are
bounded by row size eight.  Direct choice-union minimization checks 2514
extensions across the 24 classes and finds none.  This excludes $(8,8)$.

\section{The split \texorpdfstring{$(7,9)$}{(7,9)}}
\label{sec:seven-nine}

Let $|\mathcal D|=7$ and $|\mathcal E|=9$.  Injection and residual
arguments give maximum $\mathcal D$-degree three and maximum
$\mathcal E$-degree four.  In brief, degrees leaving one or two rows yield
a transversal of size at most two or at least nine private-choice blockers.
A point of $\mathcal E$-degree five leaves the unique low-count four-row
$K_4$ kernel; adding the point to zero or one of its star blockers gives,
respectively, seven or ten double-blocker rows, never nine.  A point of
$\mathcal D$-degree four would leave a three-row kernel violating (S1.1).

\subsection{A three-point row on the nine-row side}

Let $T\in\mathcal E$ have size three.  It covers all seven rows of
$\mathcal D$, so some $x\in T$ has $\mathcal D$-degree three.  The four rows
avoiding $x$ have sizes three through nine, at most nine blockers, at most
four pair blockers, and fourteen nonempty proper supports.  The 16383 active
support masks give
\[
\begin{array}{c|c|r}
\text{blockers}&\text{pair blockers}&\text{vectors}\\
\hline
7&3&1\\
8&4&48\\
9&3&12\\
9&4&75
\end{array}                                                   \tag{S9.1}
\]
or 136 vectors in eleven row-permutation classes.  Every pair blocker
acquires $x$; every other blocker may remain unchanged or acquire $x$, and
at most two further rows are added.  The extension domain retains every
admissible incidence of old points, fresh private and shared points, parallel
actual points, and all row-size-bounded multiplicities.  Independent C++ and
Python reconstructions agree on 138161 candidates and zero valid extensions.

\subsection{A three-point row only on the seven-row side}

Let $T=\{x,a,b\}\in\mathcal D$ and suppose every row of $\mathcal E$ has
size at least four.  If one point of $T$ has $\mathcal E$-degree four, the
five rows avoiding it have at most seven blockers and at most three pair
blockers.  Their point degrees are at most three.  The complete five-row
domain contains 490 vectors in eight classes, and every class has a
three-point row, contradicting the row-size assumption.

Thus all three points of $T$ have $\mathcal E$-degree three, and each of the
nine rows meets $T$ exactly once.  After deleting the three rows containing
$x$, write the remaining six as
\[
 \mathcal K=\{A_1,A_2,A_3,B_1,B_2,B_3\},
 \quad a\in A_i,\quad b\in B_i.
\]
The blocker bounds are
\[
 |\blocker(\mathcal K)|\leq7,\qquad
 |\{Q\in\blocker(\mathcal K):|Q|=2\}|\leq3.                \tag{S9.2}
\]
The three residual $A$-rows must have a unique common point $c$, and the
$B$-rows a unique common point $d$; otherwise (S1.1) gives too many total or
pair blockers.  Delete $c,d$.  Each remaining three-row residual has row
sizes two through five, no common point, and at most four blockers.  Formula
(S1.1) gives nineteen labelled multiplicity vectors in six row-permutation
classes.

A global actual point either remains on one side or matches one copy from
each side.  The degree-four bound allows exactly the matches whose local
support ranks sum to at most four.  The exhaustive coupling domain has 90898
labelled cases.  Its independently replayed pair-blocker histogram is
\[
\begin{array}{c|rrrrrr}
\text{pair blockers}&4&5&6&7&8&9\\
\hline
\text{couplings}&22485&38832&23313&5614&634&20.
\end{array}                                                   \tag{S9.3}
\]
Every coupling has at least four pair blockers, contradicting (S9.2).
Therefore neither side can contain the three-point row forced by
\cref{sec:five-eleven}; $(7,9)$ and $(9,7)$ are impossible.

\begin{theorem}[Exclusion of total sixteen]
\label{thm:middle-sixteen}
No mutual-blocker pair with minimum row size three on both sides has total
cardinality sixteen.
\end{theorem}

\begin{proof}
Blocker duality leaves the four splits handled in
\cref{sec:five-eleven,sec:six-ten,sec:eight-eight,sec:seven-nine}.
\end{proof}

\section{Why the residual searches are finite and sound}
\label{sec:finite-interface}

This section isolates the interface common to every enumeration above.
Suppose a residual clutter has $r$ labelled rows.  Map each active actual
point $u$ to its nonempty incidence support
\[
 I(u)=\{i\in[r]:u\text{ belongs to row }i\}.                \tag{S10.1}
\]
Points with the same support are not merged: their number is stored as a
positive integer multiplicity $m_I$.  Inactive points are omitted because
they occur in no row and can occur in no minimal transversal.

The reduction is finite for the following proved reasons.

\begin{enumerate}[label=(F\arabic*),leftmargin=12mm]
\item The number of possible nonempty supports is $2^r-1$; rank and
  no-full-support conclusions reduce it further in the five-row branches.
\item Every row has a proved upper bound from \cref{lem:private}.  Therefore
  every multiplicity $m_I$ is bounded by a row size containing $I$.
\item Support minimization retains only incidence masks compatible with
  clutter minimality, row-size bounds, transversal-number bounds, and the
  previously proved limits on pair blockers and point degrees.
\item A minimal transversal never uses two parallel copies of the same
  support.  For each minimal support cover $M$, the number of actual choices
  is exactly $\prod_{I\in M}m_I$.  Summing over the support covers gives the
  weighted count used in (S7.6).
\item Extension checks explicitly add all permitted incidences of existing
  points and all nonempty incidence supports of new points.  Their
  multiplicities are again bounded by row sizes.  The semantic verifier then
  expands actual labelled copies and recomputes both blockers by
  inclusion-minimal choice-union enumeration.
\end{enumerate}

Conversely, every arbitrary-finite instance maps to one enumerated
multiplicity vector by (S10.1), and expanding that vector realizes the same
row family up to inactive points and relabelling of parallel actual points.
Thus the search domain is a quotient with an explicit inverse realization,
not an assumption that $|U|$ is small.  Empty optional incidence classes,
repeated residual supports, and nonuniform rows are all present in this
description.

\subsection{Worked interface: the five-row kernel in the
\texorpdfstring{$(7,8)$}{(7,8)} split}

We trace one nontrivial branch from the natural clutter to the final Lean
contradiction.  Let $\mathcal H$ be the seven-row side of a mutual-blocker
pair, with $|\blocker(\mathcal H)|=8$.  Both sides have minimum row size three
and maximum point degree three.  Suppose a point $x$ has degree two in
$\mathcal H$, and let
\[
 \mathcal K=\{E\in\mathcal H:x\notin E\}.
 \tag{S10.2}
\]
Then $\mathcal K$ is a five-row clutter, its rows have size at least three,
and its maximum point degree is at most three.  All blocker rows in this
passage remain inclusion-minimal: whenever a pointwise minimality step is
used, \cref{lem:private} supplies a private row for that point.  Residual
blocker identities
give
\[
 \min_{Q\in\blocker(\mathcal K)}|Q|\geq2,
 \qquad |\blocker(\mathcal K)|\leq8,
 \qquad |\{Q\in\blocker(\mathcal K):|Q|=2\}|\leq3.
 \tag{S10.3}
\]
The exact five-row principle strengthens the last inequality to equality.

For $Q\in\blocker(\mathcal K)$, define its forced lift by
\[
 L_x(Q)=
 \begin{cases}
  Q\cup\{x\},&|Q|=2,\\
  Q,&|Q|>2.
 \end{cases}
 \tag{S10.4}
\]
The three pair blockers exhaust the degree available at $x$ on the eight-row
side.  Residual minimality shows that every $L_x(Q)$ belongs to
$\blocker(\mathcal H)$.  Blocker involution also gives
$\mathcal K\subseteq\blocker(\blocker(\mathcal H))$.  Therefore
$\blocker(\mathcal H)$ satisfies
\code{IsKernelCompletion K x (blocker H)}: it is an eight-row clutter with
minimum row size three and maximum degree three, contains every forced lift,
and has every row of $\mathcal K$ in its blocker.

The Lean theorem \path{G413.no_kernel_completion_of_eight} excludes exactly
this completion interface on an arbitrary finite type.  Its internal finite
case separates $|\blocker(\mathcal K)|=6,7,8$.  The first two cases use the
proved residual bound and pair-endpoint lemma.  In the eight-row case, the
three pair rows have one of the five graphs
$P_4$, $P_3\mathbin{+}K_2$, $3K_2$, $K_3$, or $K_{1,3}$; the declarations
assembled in \path{G413.fiveRow_eight_blocker_impossible} exclude all five.
The $3K_2$ branch uses the support quotient (S10.1) together with a proved
lift back to arbitrary actual points, so parallel points are not identified.
More explicitly, its canonical finite domain is a triple of
\code{ValidFiveSupportPair} values (50 choices per pair) satisfying
\code{ThreeK2FiveSupportsCompatible}; \code{threeK2FiniteGood} is the
enumerated predicate.  The theorem
\path{G413.threeK2_finite_support_classification} proves the finite split, while
\path{G413.threeK2_actual_support_classification} constructs those three support
pairs from arbitrary actual rows and transfers the result.  The decoder and
certificate semantics are recorded in the following declarations:
\begin{center}
\small
\begin{tabular}{l}
\path{G413.threeK2Decode},\\
\path{G413.finitePrivateBlockerCertificate_eq_true_iff},\\
\path{G413.finiteBlocker_eq_blocker}.
\end{tabular}
\end{center}
The arbitrary-finite contradiction is
\path{G413.threeK2_pairGraph_impossible}.

The fully qualified name of the final declaration is
\begin{center}
\small
\path{MiddleSevenEightFiveKernelBranch.}\allowbreak%
\path{sevenRows_degreeTwo_impossible}.
\end{center}
This declaration constructs the interface (S10.2)--(S10.4) from the natural
seven-row family and applies
\path{G413.no_kernel_completion_of_eight}.  Thus this branch does not
deduce an arbitrary-ground conclusion from the executable counts.  The
programs replay the finite support classification; the Lean bridge proves
that every arbitrary-finite realization enters it and that every decoded
case has the required semantics.

The executable sources and their saved outputs are organized as follows.
Every displayed path is relative to
\path{Problems/P0054/experiments/} inside the artifact root.

\begin{longtable}{>{\raggedright\arraybackslash}p{0.19\textwidth}
                  >{\raggedright\arraybackslash}p{0.35\textwidth}
                  >{\raggedright\arraybackslash}p{0.36\textwidth}}
\toprule
Claim & Primary enumeration & Independent replay or semantic check \\
\midrule
\endfirsthead
\toprule
Claim & Primary enumeration & Independent replay or semantic check \\
\midrule
\endhead
Fano adjacency
& \path{g4.12-exact-fifty-four/fano-adjacency.cpp}
& \path{g4.12-exact-fifty-four/verify.py} \\
$(6,10)$, degree three
& \path{g4.12-exact-fifty-four/middle-six-ten.cpp}
& \path{g4.12-exact-fifty-four/verify-middle-six-ten.py} \\
$(6,10)$, low degree
& \path{g4.12-exact-fifty-four/middle-six-ten-low-degree.cpp}
& \path{g4.12-exact-fifty-four/middle-six-ten-low-degree.py} \\
$(8,8)$, four rows
& \path{g4.12-exact-fifty-four/middle-four-row-kernels.py}
& saved JSON is rechecked by the extension verifier \\
$(8,8)$, five rows
& \path{g4.12-exact-fifty-four/middle-eight-eight-five-row.cpp}
& \path{g4.12-exact-fifty-four/middle-eight-eight-extension.py} \\
$(7,9)$, four rows
& \path{g4.12-exact-fifty-four/middle-seven-nine-four-row.cpp}
& \path{g4.12-exact-fifty-four/middle-four-row-kernels.py} \\
$(7,9)$, extensions
& \path{g4.12-exact-fifty-four/middle-seven-nine-extension.cpp}
& \path{g4.12-exact-fifty-four/middle-seven-nine-extension.py} \\
$(7,9)$, couplings
& \path{g4.12-exact-fifty-four/middle-seven-nine-six-row.cpp}
& \path{g4.12-exact-fifty-four/middle-seven-nine-six-row.py} \\
52-profile kernels
& \path{g4.10-coupled-fifty-two/four-row-kernel.cpp}
& \path{g4.10-coupled-fifty-two/support-multigraph-checker.py} \\
53-profile kernels
& \path{g4.11-coupled-fifty-three/five-row-kernel.cpp}
& \path{g4.11-coupled-fifty-three/verify.py} \\
\bottomrule
\end{longtable}

Each experiment directory contains a manifest with toolchain information and
SHA-256 digests.  The saved numeric output is \computed.  The conclusion
about arbitrary finite clutters follows only after the proved interface
(F1)--(F5); in the final artifact this interface and the classifications are
also checked in Lean.

\clearpage
\begin{landscape}
\section{Theorem--proof--program--Lean crosswalk}
\label{sec:crosswalk}

The table below gives the shortest review path.  Every P0054 name is complete
after the common prefix \code{AiMathLab.P0054.}  For the first infrastructure
entry, prepend \code{AiMathLab} and a period instead.  Names never inherit a
namespace from the preceding entry.  Source files are at the repository root.

\small
\begin{longtable}{>{\raggedright\arraybackslash}p{0.18\textwidth}
                  >{\raggedright\arraybackslash}p{0.24\textwidth}
                  >{\raggedright\arraybackslash}p{0.49\textwidth}}
\toprule
Natural statement & Human proof / executable evidence & Lean declaration \\
\midrule
\endfirsthead
\toprule
Natural statement & Human proof / executable evidence & Lean declaration \\
\midrule
\endhead
Homogeneous atom from signatures
& Main theorem ``Finite canonical reduction''
& \path{SaturatedSperner.HomogeneousAtomExistence.exists_homogeneousAtom_of_signature_bound} \\
Canonical trace blocker equality and row bounds
& Main canonical-reduction theorem
& \path{CanonicalProfiles.canonicalLayer_mutual_blockers};
  \path{CanonicalProfiles.smallTrace_rows_cardAtLeast};
  \path{CanonicalProfiles.largeComplementTrace_rows_cardAtLeast} \\
Layer trace cardinality and total sum
& Main equations (CR3) and the canonical partition
& \path{CanonicalProfiles.canonicalLayer_trace_card_sum};
  \path{CanonicalProfiles.sum_canonicalLayer_cards_eq} \\
$m(2,3)=9$
& Main proof of the local blocker theorem
& \path{Sat6.local_lower}, together with
  \path{Sat6.witness_blocker_large}, \path{Sat6.witness_blocker_small}, and
  \path{Sat6.witness_total_card} \\
$m(2,4)=12$
& Main theorem and its product-bound appendix
& \path{AdjacentExact.mutual_blocker_total_card_ge_twelve};
  \path{AdjacentExact.witness_blocker_two};
  \path{AdjacentExact.witness_blocker_four};
  \path{AdjacentExact.adjacent_local_parameter_exact} \\
Parameterized 30-member construction
& Main six-layer upper bound
& \path{Sat6Parameterized30.family_card};
  \path{Sat6Parameterized30.family_saturated};
  \path{Sat6Parameterized30.full_certificate} \\
Finite-ground lower bound 30
& Canonical reduction plus $m(2,3)=9$
& \path{Sat6Exact.family_card_lower} \\
$\sat(6)=30$
& Main exact theorem
& \path{Sat6StableExact.sat_six_eq_thirty} \\
$m(3,3)\geq14$
& \Cref{sec:m33-lower}
& \path{MiddleLower.mutual_blocker_total_card_ge_fourteen} \\
Equality 14 is Fano
& \Cref{sec:m33-classification}; G4.10 replay
& \path{MiddleFourteenFano.middle_fourteen_is_fano} \\
No middle total 15
& \Cref{sec:m33-classification}; G4.11 replay
& \path{MiddleFifteen.middle_fifteen_impossible} \\
Five-row completion in the $(7,8)$ split
& \Cref{sec:finite-interface}, worked interface
& \path{G413.no_kernel_completion_of_eight};
  \path{MiddleSevenEightFiveKernelBranch.sevenRows_degreeTwo_impossible} \\
$(5,11)$ endpoint
& \Cref{sec:five-eleven}
& \path{MiddleFive.blocker_card_ge_twelve} \\
$(6,10)$ impossible
& \Cref{sec:six-ten}; two independent replays
& \path{G419.sixTen_impossible} \\
$(8,8)$ impossible
& \Cref{sec:eight-eight}; four- and five-row replays
& \path{G420.eightEight_impossible} \\
$(7,9)$ impossible
& \Cref{sec:seven-nine}; extension and coupling replays
& \path{MiddleSevenNineClosed.sevenNine_impossible} \\
No middle total 16
& \Cref{thm:middle-sixteen}
& \path{MiddleSixteenClosed.middle_sixteen_impossible} \\
Fano $(7,6)$ and $(8,5)$
& \Cref{sec:fano-profile}; Fano adjacency replay
& \path{FanoAdjacentSevenSixActual.fano_seven_six_impossible};
  \path{FanoAdjacentEightFiveActual.fano_eight_five_impossible} \\
Fano $(9,4)$
& \Cref{lem:four-ten}
& \path{AdjacentFourIntersecting.four_pairwise_intersecting_blocker_card_ge_ten} \\
All Fano-adjacent splits
& \Cref{sec:fano-profile}
& \path{FanoAdjacentClosed.fano_adjacent_impossible} \\
Canonical lower bridge
& Main profiles and both exact exclusions
& \path{SevenFanoBridge.family_card_lower_of_fano_adjacent} \\
Parameterized 55-member construction
& Main construction table and saturation proof
& \path{Parameterized55.family_card};
  \path{Parameterized55.family_saturated};
  \path{Parameterized55.full_certificate} \\
$\sat(7)=55$
& Main exact theorem
& \path{Sat7StableExact.sat_seven_eq_fifty_five} \\
\bottomrule
\end{longtable}
\end{landscape}
\clearpage

The stable declarations have the proposition types
\[
 \code{IsStableSaturationNumber 6 30},\qquad
 \code{IsStableSaturationNumber 7 55}.
\]
This predicate asserts the existence of a threshold beyond which both an
attaining saturated family and a universal lower bound hold on every
labelled finite ground set.  The explicit witness thresholds in the two Lean
proofs are $3\cdot2^{30}$ and $3\cdot2^{55}$.  These are existence thresholds,
not claims of optimal stabilization.

\section{Reproduction and trust boundary}
\label{sec:reproduction}

The versioned artifact is available from
\href{https://github.com/lailai0916/saturated-sperner-6-7/releases/tag/v1.0.0}
{GitHub release v1.0.0} and
\href{https://doi.org/10.5281/zenodo.21679078}{Zenodo DOI
10.5281/zenodo.21679078}.  That immutable record is the public baseline and
predates the present revision.  The revised sources and PDFs must be deposited
as a new GitHub/Zenodo version rather than replacing v1.0.0; until its version
DOI is inserted into the submission files, the public artifact and this
revision are not version-synchronized.

From the artifact root, the two principal Lean checks are
\begin{verbatim}
lake build AiMathLab AiMathLab.P0054Sat7StableExact
lake env lean Problems/P0054/formal/Main.lean
\end{verbatim}
The complete recorded build checks 17488 jobs.  The formal entry point prints
the axioms of the two final declarations; the accepted list is exactly
\code{propext}, \code{Classical.choice}, and \code{Quot.sound}.  A source
scan rejects \code{sorry}, \code{admit}, \code{axiom}, \code{unsafe},
\code{native\_decide}, and \code{run\_tac}.

The explicit 55-member construction is checked independently in three ways:
direct expansion on the eleven-point realization, a semantic Python
verifier, and a Lean certificate.  In Lean, all $2^{11}=2048$ subsets are
covered by a frozen saturation table; the 1993 external subsets carry an
eight-member strict chain after insertion.  The generator is untrusted:
Lean checks the length, membership, strict containment, and occurrence
conditions of every emitted row.

The status boundary is therefore:
\begin{itemize}[leftmargin=8mm]
\item the executable tables and counts in this supplement are \computed;
\item the self-contained arguments displayed in full are \proved;
\item compressed finite-classification narratives are audit maps, not
  standalone arbitrary-finite proofs;
\item the final stable values and all classification dependencies are
  \mbox{\formalized};
\item novelty, priority, editorial suitability, and any claim beyond the two
  exact stable values are not mathematical consequences of these checks.
\end{itemize}

\end{document}
